\section{\texttt{XAppCommitment}: staking without a P fund ledger}
\label{sec:xappcommitment}

The XLock (Section~\ref{sec:xlock}) is the universal claim handle for a
\emph{transient} commitment --- an order, a bridge transfer --- that
settles and releases on a short horizon. Staking is the same idea on a
long horizon, and it is where the parallel-ledger problem is sharpest:
the P-Chain today owns a second spendable UTXO set and \emph{mints new
value} on its own state. This section defines the object that removes
that --- \textbf{\texttt{XAppCommitment}} --- and its staking
specialization \textbf{\texttt{XStakeCommitment}}.

The thesis of this section is narrow and load-bearing, so we state it
before the mechanism: \textbf{Lux does not require economic slashing for
staking safety.} Stake safety is asset exclusivity --- the committed
principal can be used in exactly one place --- and exclusivity is a
property of X's single consume-once state machine, not of a threat to
burn the principal. Misbehavior is answered by withholding the reward,
ejecting the validator, and revoking its capability NFT; the principal
is always returned.

\subsection{The key idea: P owns no spendable fund ledger}

Under XSS, \textbf{P owns no spendable fund ledger.} X finalizes a
\emph{commitment} of funds to P, and P treats that commitment as stake
weight --- nothing more. P never holds the principal, never returns it,
and never mints the reward; it \emph{computes} the validator lifecycle
and asks X to settle. The funds stay locked on X for the entire bond and
move only by an X transition.

This is the same move XSS makes for the DEX, generalized. A DEX order
locks value on X and the matcher holds a \emph{claim}; a validator
commitment locks value on X and the validator set holds a \emph{claim}.
One mechanism, two purposes.

We say P owns no spendable fund ledger, not that P owns no ledger at all.
P keeps lifecycle and reward-eligibility bookkeeping --- a delegatee
reward accumulator, uptime, the validator set --- and that bookkeeping
is real (the delegatee accumulator is \texttt{GetDelegateeReward} /
\texttt{AddRewardUTXO} in \texttt{vms/platformvm/txs/executor/\allowbreak
proposal\_tx\_executor.go:327--359}, \IMPL{}). The precise claim is that
none of P's bookkeeping is \emph{spendable}: it is eligibility data X
consumes when it settles, never an output a P transaction can move.

\subsection{The object}

\texttt{XAppCommitment} generalizes \texttt{XLock} so a single object
covers stake, trade, and bridge commitments to any app-chain module:

\begin{lstlisting}
XAppCommitment {
  commitmentID    // unique id, derived from the committing X tx
  sourceUTXOIDs   // the spendable X UTXOs consumed to fund it
  owner           // who the principal returns to (release/abort target)
  assetID         // single asset; one commitment locks one asset
  amount          // == sum(value(sourceUTXOIDs))
  appChainID      // which app-chain may reference this commitment
  appModuleID     // which module on that chain (e.g. validator_set)
  purpose         // stake | trade | bridge
  constraints     // interval [start,end], minAmount, ...
  activated       // X-visible flag; set by AcceptStake, never unset
  consumed        // set true exactly once, by a settlement transition
}
\end{lstlisting}

\noindent For staking the instantiation is fixed: \texttt{purpose=stake},
\texttt{appChainID=P}, \texttt{appModuleID=validator\_set}, and
\texttt{constraints} carries the bond interval and the minimum bond. We
call that instantiation \texttt{XStakeCommitment}; it is not a new type,
it is \texttt{XAppCommitment} with these fields pinned. Note what is
\emph{absent}: there is no \texttt{slashable} flag and no slashed-fraction
field. The object cannot express confiscation of principal because the
protocol never confiscates principal (Theorem~\ref{thm:noslash}).

\DESIGN{} No such object exists today. Its relationship to the existing
code is precisely the relationship \texttt{XLock} has to the P-Chain's
\texttt{StakeOuts}: \texttt{XStakeCommitment} is what \texttt{StakeOuts}
becomes when the lock moves off P and onto X, and the principal stops
being a P-Chain output (the \texttt{StakeOuts} citation and the full
contradiction are in Section~\ref{sec:stake-contradiction} below).

\subsection{The transition lattice: exactly six transitions}
\label{sec:stake-transitions}

A commitment moves through a fixed state machine. The six transitions
below are the \emph{only} ways its \texttt{consumed} or \texttt{activated}
flag is touched or its funds move; X rejects anything else. There is no
burn transition and no confiscation transition --- those are not omitted
for brevity, they are absent by design.

\begin{enumerate}[leftmargin=2.4em,label=\textbf{T\arabic*.}]
\item \textbf{\texttt{CommitStake}} \;(X UTXO $\to$
  \texttt{XStakeCommitment}, \texttt{target=P}).
  Consume \texttt{sourceUTXOIDs}; create one \texttt{XStakeCommitment}
  with \texttt{activated=false}, \texttt{consumed=false}. Conservation:
  net zero --- spendable value enters the locked pool, total unchanged.

\item \textbf{\texttt{AcceptStake}} \;(P references the finalized
  commitment). P's \texttt{validator\_set} module records a validator
  claim keyed by \texttt{commitmentID}; X sets \texttt{activated=true}.
  \textbf{This touches no X value and does not set \texttt{consumed}} ---
  it is a pure P-side computation over a finalized X fact, plus the one
  X-visible flag flip that closes the activate/abort race
  (Section~\ref{sec:activate-race}). Capability-gated: a
  validator-permission NFT (\texttt{RequireNFT},
  Section~\ref{sec:primitives}) may gate acceptance; absent it,
  acceptance fails closed.

\item \textbf{\texttt{CompleteStake}} \;(P completion proof $\to$ X
  releases principal $+$ bounded reward if eligible). On bond expiry
  with the validator having met the reward conditions, P emits a
  completion certificate; X consumes the commitment
  (\texttt{consumed=true}) and creates (a)~the principal output back to
  \texttt{owner} and (b)~a reward output, where the reward amount is
  \textbf{recomputed by X} from its own schedule
  (Section~\ref{sec:rewardmodel}), not taken on P's word. This is the
  only transition that may increase $\mathrm{total}(a)$
  (Section~\ref{sec:rewardinflation}).

\item \textbf{\texttt{AbortStake}} \;(stake never accepted / invalid /
  expired $\to$ X releases principal \emph{only}). If the commitment was
  never accepted (\texttt{activated=false}) and is invalid or past
  \texttt{constraints.end}, X consumes the commitment and returns the
  full \texttt{amount} to \texttt{owner}. No reward. This is the staking
  instance of the XLock liveness guarantee
  (Section~\ref{sec:xlock}): value can always come home if P fails to
  accept. Conservation: net zero.

\item \textbf{\texttt{FailReward}} \;(validator accepted but did not meet
  reward conditions $\to$ X releases principal \emph{only}). The
  commitment was accepted (\texttt{activated=true}) but the validator was
  offline or insufficiently correct, so P's outcome proof says
  ``no reward.'' X consumes the commitment and returns the full
  \texttt{amount} to \texttt{owner}. \textbf{The principal is returned in
  full; only the reward is withheld.} Conservation: net zero. This is the
  no-slash analog of the legacy abort-block path: same economics
  (principal back, no reward), no burn.

\item \textbf{\texttt{RevokeOperator}} \;(governance / authorization
  event revokes a validator or plugin-operator capability). An authorized
  revocation (e.g. via a \texttt{DisableAuth}-style authorization, modeled
  on \texttt{vms/platformvm/txs/\allowbreak disable\_l1\_validator\_tx.go:15--22},
  \IMPL{}) ejects the operator and revokes its capability NFT. It
  \textbf{does not consume the stake commitment and burns nothing}; the
  staked principal settles through \texttt{CompleteStake} or
  \texttt{FailReward} on the normal schedule (reward only if the
  pre-revocation interval qualified). Revocation removes future
  capability, not past principal.
\end{enumerate}

\noindent \texttt{CompleteStake}, \texttt{AbortStake}, and
\texttt{FailReward} are mutually exclusive and each sets
\texttt{consumed} exactly once; \texttt{CommitStake} and
\texttt{AcceptStake} precede them; \texttt{RevokeOperator} is orthogonal
(it sets neither \texttt{consumed} nor \texttt{activated} and may fire at
any point a capability exists). There is no transition that makes the
committed principal spendable while the validator is active, and no
transition that destroys it.

\paragraph{Economic outcomes, enumerated.} Every terminal state returns
principal in full. The reward and capability differ:
\begin{center}
\begin{tabular}{@{}lll@{}}
\toprule
Outcome & Principal & Reward / capability \\
\midrule
success (\texttt{CompleteStake})        & returned & $+$ bounded reward \\
offline / incorrect (\texttt{FailReward}) & returned & zero reward \\
never accepted (\texttt{AbortStake})     & returned & zero reward \\
misconduct / authz (\texttt{RevokeOperator}) & returned & capability revoked, ejected \\
\bottomrule
\end{tabular}
\end{center}
\noindent \textbf{Principal is never lost} in any column.

\subsection{The activate/abort race}
\label{sec:activate-race}

\texttt{AcceptStake} and \texttt{AbortStake}/Expire must not both fire on
the same commitment --- one would mint a validator claim while the other
returns the bond. The \texttt{activated} flag serializes them on X:
\texttt{AcceptStake} sets \texttt{activated=true} and is the only writer
of that flag; \texttt{AbortStake}/Expire is admissible \emph{only if}
\texttt{activated=false} at the height it executes. Because both observe
a single finalized X total order (A1, Section~\ref{sec:stake-theorems}),
exactly one wins: if acceptance finalized first, the abort is rejected
(\texttt{activated} is set); if the expiry finalized first, acceptance is
rejected (the commitment is consumed/expired). \DESIGN{} The flag is the
new part; the serialization it relies on is X's existing single-order
finality.

\subsection{P's validation duty is a predicate, not custody}

Under this model P validates a single predicate before
\texttt{AcceptStake} and before honoring any reference to the commitment:

\begin{quote}
\emph{A finalized X commitment exists, targeted to P
(\texttt{appChainID=P}, \texttt{appModuleID=validator\_set}), for this
owner / validator / amount / asset / interval, that is not expired and
not consumed.}
\end{quote}

\noindent That is the whole of P's coupling to value. P does not hold
the funds, does not track a spendable balance, and does not decide when
value becomes spendable --- it checks a T2 fact (a finalized X
commitment) and runs its T1/T3-free computation (the validator
lifecycle and reward eligibility) on top of it. This is the access model
of Section~\ref{sec:mesh} applied to staking: P is a computing module
over X-committed value.

\subsection{Anti-double-spend: exclusivity without P custody}
\label{sec:stake-exclusivity}

The property that makes this safe \emph{without} P holding custody:

\begin{quote}
\textbf{Commitment exclusivity.} Once \texttt{CommitStake} consumes
\texttt{sourceUTXOIDs}, that value cannot be traded on D, bridged on B,
spent on X, shielded on Z, or re-staked --- because the originating X
UTXOs are \emph{consumed}, and a consumed UTXO has no spend path on any
chain.
\end{quote}

\noindent The argument is the consume-once discipline of the UTXO
engine, not a new lock service. \texttt{CommitStake} consumes its
source UTXOs exactly as any X transaction consumes its inputs (the
\IMPL{} consume-once on inputs, \texttt{vms/xvm/txs/executor/\allowbreak
semantic\_verifier.go:152--187} \texttt{verifyTransfer}; the same
consume-once the atomic mailbox relies on,
\texttt{lux/chains/dexvm/atomic.go:126--138}). After consumption the
value exists in exactly one place --- inside the
\texttt{XStakeCommitment} on X --- and is reachable only by the
settlement transitions above, each of which sets \texttt{consumed} once.
Exclusivity across \emph{all} uses (trade / bridge / spend / shield /
re-stake) is therefore free: there is one X state per unit
(Section~\ref{sec:model}), and a committed unit is in the
\emph{locked} state under one named commitment. No P-side custody, no
second ledger, no reconciliation --- the exclusion is a property of X's
single state machine.

Crucially, this exclusivity \emph{rides on} X's consensus safety; it does
not manufacture it. ``Consume-once'' is meaningful only because X
provides a single finalized total order in which each source UTXO is
spent at most once (A1 below). Atomicity of the
consume-and-commit step prevents double-\emph{use} of a finalized UTXO;
it does not by itself guarantee that the order is final --- that is the
job of consensus. The staking theorems are therefore stated as
implications from consensus safety, never as replacements for it.

\subsection{Staking safety theorems}
\label{sec:stake-theorems}

We make the kernel above precise. The theorems are conditional on
consensus safety and reduce to the ledger-level consume-once /
conservation results proved in the XSS soundness note
(\texttt{lux/proofs/xss-double-spend-inflation-soundness.tex}); they
restate those results in staking terms and add the no-slash corollary.

\paragraph{A1 (consensus assumption).} X provides a single finalized
total order of accepted transactions, and atomically applies, per
transaction, the consumption of its input UTXOs together with the
creation of its outputs and commitments. (This is X's Quasar finality
guarantee; we assume it, we do not prove it here.)

\begin{theorem}[No double-spend of committed stake]
\label{thm:nodouble}
Given A1, once an X UTXO $u$ is consumed by a \texttt{CommitStake} into an
\texttt{XStakeCommitment} (\texttt{purpose=stake}, \texttt{target=P}), no
other valid X transition can consume $u$.
\end{theorem}
\begin{proof}[Proof sketch]
By A1 the consumption of $u$ is applied atomically with the creation of
the commitment in one finalized transaction, and $u$ is thereafter
absent from the unspent set. Any later transaction naming $u$ as an input
fails the consume-once check (\IMPL{}
\texttt{semantic\_verifier.go:152--187}); the single total order admits no
``later'' that precedes the consumption. Hence $u$ is spent exactly once.
The conclusion is the staking instance of the ledger consume-once lemma
in the soundness note. Note the dependence on A1: without finality of the
order, two branches could each consume $u$ --- atomicity alone does not
exclude that, consensus does.
\end{proof}

\begin{theorem}[No inflation from staking settlement]
\label{thm:noinfl}
Given A1 and an X-enforced reward schedule
(Section~\ref{sec:rewardmodel}), staking settlement creates only
(1)~returned principal already committed in the consumed
\texttt{XStakeCommitment}, and (2)~a bounded protocol reward permitted by
that schedule. P cannot mint: P produces only lifecycle evidence
(completion / failure certificates), which X consumes; the spendable
outputs are created by X.
\end{theorem}
\begin{proof}[Proof sketch]
\texttt{AbortStake} and \texttt{FailReward} create exactly the principal
recorded in \texttt{amount} and nothing else (net-zero). \texttt{CompleteStake}
creates that same principal plus a reward output whose amount X recomputes
from \texttt{XRewardSchedule}; by construction the schedule bounds the
reward (the legacy bound is the supply invariant
\texttt{state\_changes.go:173--175}: \texttt{rewards.Calculate} can never
push \texttt{supply + reward} past \texttt{maximumSupply}). P emits no
output and holds no mint authority: \texttt{RewardValidatorTx} has empty
\texttt{InputIDs}/\texttt{Outputs}
(\texttt{vms/platformvm/txs/reward\_validator\_tx.go:45--50}, \IMPL{}). Thus the
only new value is the schedule-bounded reward, which reduces to the
inflation-soundness result of the proofs note.
\end{proof}

\begin{theorem}[No slashing required for stake safety]
\label{thm:noslash}
Stake safety (Theorems~\ref{thm:nodouble},~\ref{thm:noinfl}) does not
depend on burning, confiscating, or redirecting committed principal.
Misbehavior affects reward eligibility (\texttt{FailReward}) or
authorization status (\texttt{RevokeOperator}); neither alters the
conservation of committed principal, which is returned in full on every
terminal transition.
\end{theorem}
\begin{proof}[Proof sketch]
Exclusivity (Theorem~\ref{thm:nodouble}) follows from consume-once under
A1 and is independent of any burn: no transition in
Section~\ref{sec:stake-transitions} reduces \texttt{amount} below the
committed principal, and the only supply-increasing transition is the
bounded reward (Theorem~\ref{thm:noinfl}). The three terminal transitions
(\texttt{CompleteStake}, \texttt{AbortStake}, \texttt{FailReward}) each
return the full principal; \texttt{RevokeOperator} consumes nothing.
Therefore the safety properties hold with the slash transition removed,
which establishes the claim. (Liveness incentives are a separate concern,
handled by reward non-payment and ejection, not by conservation.)
\end{proof}

\subsection{The reward model: X recomputes, P proves the branch}
\label{sec:rewardmodel}

The reward is a function evaluated \emph{by X}:
\[
\mathrm{rewardAmount} = \texttt{XRewardSchedule}\bigl(
  \text{validatorID},\,\text{stakeCommitmentID},\,\text{amount},
\]
\[
  \text{startTime},\,\text{endTime},\,\text{uptime/correctness proof},
  \,\text{network params}\bigr).
\]
\texttt{CompleteStake} consumes the commitment and creates a principal
output plus a reward output of \texttt{rewardAmount} when eligible;
\texttt{FailReward} consumes the commitment and creates the principal
output only. No P/D/B/T transition mints: there is no
\texttt{Fx.CreateOutput} of a reward on any chain but X.

\paragraph{C2 --- the commit/abort fork is a BFT vote, not an
attestation.} What P hands X is \emph{not} a signed ``proof of good
standing.'' The legacy reward path settles through a P-Chain
\emph{proposal vote}: \texttt{RewardValidatorTx} carries no credentials
at all --- its struct is just a \texttt{TxID} plus cached bytes
(\texttt{reward\_validator\_tx.go:28--32}), and it declares empty inputs
and outputs (\texttt{:45--50}). The reward-vs-no-reward decision is made
by the option-block oracle: \texttt{prefersCommit} fetches the staker,
computes a local \texttt{float64} uptime via
\texttt{CalculateUptimePercentFrom}, and prefers the \texttt{CommitBlock}
iff \texttt{uptime >= expectedUptimePercentage}
(\texttt{vms/platformvm/block/executor/options.go:102--148}); on error it
falls back to commit to err toward over-rewarding (\texttt{:79--85}).
Validators then vote a \texttt{CommitBlock} or \texttt{AbortBlock} and
BFT finalizes one branch.

\CRYPTO{}\,/\,\DESIGN{} The X-side conditioning object for
\texttt{CompleteStake}/\texttt{FailReward} is therefore a \emph{finality
proof of which block branch P finalized} (CommitBlock $\Rightarrow$
reward eligible, AbortBlock $\Rightarrow$ no reward), \emph{not} an
attestation of good standing. X reads the branch P committed and
\textbf{recomputes the reward itself} from \texttt{XRewardSchedule}; it
never trusts a P-supplied amount. This is what keeps
Theorem~\ref{thm:noinfl} honest: a compromised P can at most steer the
branch (a liveness/eligibility lie), it cannot inflate, because X owns
the schedule. The exact branch-finality proof and the schedule are
cryptographer- and economics-review items
(Section~\ref{sec:open}).

\paragraph{C1 --- supply-timing reconciliation (the one explicit
reward-schedule design decision).} This target model mints the reward at
\texttt{CompleteStake}. The current code does the opposite: it
\textbf{pre-mints at activation}. When a pending staker is promoted to
current, \texttt{state\_changes.go} computes \texttt{potentialReward} and
immediately does
\texttt{SetCurrentSupply(supply + potentialReward)} before
\texttt{PutCurrentValidator}
(\texttt{vms/platformvm/txs/executor/state\_changes.go:166--175}, \IMPL{}),
then claws the reward back on the abort branch
(\texttt{proposal\_tx\_executor.go:294--295} write the refund to
\emph{both} commit and abort states, while the reward UTXO is added only
on commit, \texttt{:300--322}). \texttt{XRewardSchedule} can preserve the
\emph{per-validator} amount under either timing, but \textbf{mint-at-completion
versus mint-at-admission changes how supply compounds}: pre-minting
inflates \texttt{currentSupply} (and thus the denominator
\texttt{rewards.Calculate} sees for the next staker) at admission rather
than at completion. We default to the spec --- \textbf{mint at
\texttt{CompleteStake}} --- and flag the supply-timing as the single
explicit \DESIGN{} decision the reward-schedule design must ratify in P4.

\subsection{Delegation and L1 continuous-fee staking}
\label{sec:stake-delegation}

Two existing staking shapes must map onto the lattice; both already
exist in code and neither requires slashing.

\paragraph{Delegation.} A delegator's stake is a second
\texttt{XStakeCommitment} that shares a validator's weight ceiling.
Admissibility is cross-commitment: the combined weight on a validator may
not exceed the limit (\texttt{overDelegated}/\texttt{GetMaxWeight},
\texttt{vms/platformvm/txs/executor/\allowbreak staker\_tx\_verification\_helpers.go:116--133},
\IMPL{}), and the delegation interval must be a subset of the validator's
(\texttt{BoundedBy}). The delegatee's cut accrues in a reward accumulator
that X reads at \texttt{CompleteStake}
(\texttt{GetDelegateeReward} $\to$ \texttt{AddRewardUTXO},
\texttt{proposal\_tx\_executor.go:327--359}, \IMPL{}). This is the precise
sense in which ``P holds no \emph{spendable} ledger'' is the right
phrasing: P holds an eligibility accumulator, which X turns into a
reward output --- P never holds a spendable balance.

\paragraph{L1 continuous-fee staking.} Sovereign-L1 validators do not
post a fixed bond for a fixed interval; they hold a continuously-debited
balance. \texttt{IncreaseL1ValidatorBalanceTx} tops the balance up by at most the
net \$LUX it consumes (the \texttt{Balance} field is bounded by
inputs less outputs less fee,
\texttt{vms/platformvm/txs/\allowbreak increase\_l1\_validator\_balance\_tx.go:20--26},
\IMPL{}) and \texttt{DisableL1ValidatorTx} is an authorized early exit
(\texttt{DisableAuth.Verify()},
\texttt{vms/platformvm/txs/\allowbreak disable\_l1\_validator\_tx.go:15--22}, \IMPL{})
--- neither burns. \DESIGN{} The lattice covers this with a
\emph{continuous-balance variant} of \texttt{XStakeCommitment}: top-up is
a \texttt{CommitStake} that augments an existing commitment's
\texttt{amount}; early-exit is an authorized \texttt{AbortStake} that
returns the remaining balance. Where the continuous variant is not built,
the base lattice is explicitly scoped to fixed-bond staking and L1
continuous staking remains on the legacy path until P4 --- we say so
rather than pretend the fixed-bond lattice already models it.

\paragraph{The complete outcome lattice (no slash).} Collecting the
above, a validator's economic outcome is one of:
\[
\{\,\text{activate-fail} \to \text{return},\;
   \text{under-uptime}/\texttt{FailReward} \to \text{return, no reward},
\]
\[
   \text{good-standing}/\texttt{CompleteStake} \to \text{return} + \text{reward},\;
   \text{L1-topup},\; \text{L1-early-exit}\,\}.
\]
No element burns principal.

\subsection{Operator model and accountability without burn}
\label{sec:operator-accountability}

The same no-slash discipline governs the non-staking operators ---
DEX matchers and bridge signers --- that settle through X. An operator's
authority is a capability NFT (Section~\ref{sec:primitives}); its
accountability is the NFT's revocability, not a bond at risk.

\begin{itemize}[leftmargin=2em,nosep]
\item \textbf{Invalid receipt:} the settlement is \emph{rejected} at X's
  \texttt{VerifyChainCommitment} --- no state moves, nothing to slash.
\item \textbf{Repeated bad output / equivocation:} the operator's
  capability NFT is \emph{revoked}; it can no longer activate the module.
\item \textbf{Unavailable:} the operator earns no fees or rewards and is
  deprioritized / its endpoint marked unhealthy --- non-payment, not
  confiscation.
\item \textbf{Malicious:} \emph{ejection} and NFT revocation.
\end{itemize}

\noindent To make equivocation accountable we record evidence, not a
penalty:
\begin{lstlisting}
EquivocationEvidence {
  operatorID            // the capability-NFT holder
  chainID               // the module/chain it equivocated on
  conflictingMessages   // the two (or more) conflicting signed outputs
  signatures            // operator signatures binding it to each
  height                // X height at which the conflict is recorded
}
\end{lstlisting}
\noindent \textbf{Allowed consequences} of recorded evidence: revoke the
capability NFT, eject the operator, mark its endpoint unhealthy, deny
future rewards, pause the module's route. \textbf{Disallowed in the base
model}: burn principal, confiscate stake, redirect stake to a treasury.
Equivocation evidence is for accountability and capability revocation; it
does not imply stake confiscation. \DESIGN{} The evidence object and its
consequence set are new; the capability NFTs it acts on are the \IMPL{}
\texttt{nftfx}/\texttt{propertyfx} outputs (\texttt{vms/xvm/fxs/fx.go}).

\subsection{The slashing-policy statement}
\label{sec:slash-policy}

We state the policy verbatim, as the authoritative position of this LP:

\begin{quote}
Lux does not require economic slashing for staking or plugin-chain
safety. \texttt{XAppCommitment}s provide asset exclusivity by consuming
source X UTXOs into purpose-bound commitments. P manages validator
lifecycle and reward eligibility, but cannot mint, burn, or confiscate
stake. X settles principal and rewards according to X-level rules.
Misbehavior is handled by non-payment of rewards, validator/operator
ejection, and capability NFT revocation. Equivocation evidence may be
recorded for accountability, but it does not imply stake confiscation in
the base protocol.
\end{quote}

\subsection{The contradiction in current code}
\label{sec:stake-contradiction}

\texttt{XStakeCommitment} is a migration target precisely because the
P-Chain today does the opposite on three counts. All three are
\IMPL{} (they are what the code does now) and all three contradict
``other chains compute, X settles'':

\begin{enumerate}[leftmargin=2em]
\item \textbf{P keeps its own spendable UTXO ledger.} The platform state
  holds \texttt{modifiedUTXOs}, \texttt{utxoDB}, and \texttt{utxoState}
  --- a second, independent UTXO set
  (\texttt{vms/platformvm/state/state.go:429--431}). \IMPL{}
\item \textbf{P holds the principal as its own locked output.} Stake is
  \texttt{StakeOuts}, P-Chain \texttt{TransferableOutput}s held on P,
  not a lock on X
  (\texttt{vms/platformvm/txs/\allowbreak add\_permissionless\_validator\_tx.go:49--50}).
  \IMPL{}
\item \textbf{P mints reward value and returns principal on its own
  state.} At end-of-stake the proposal executor refunds the stake by
  adding UTXOs to P's state on both forks (\texttt{onCommitState.AddUTXO}
  / \texttt{onAbortState.AddUTXO},
  \texttt{vms/platformvm/txs/executor/\allowbreak
  proposal\_tx\_executor.go:294--295}) \emph{and mints the reward} on the
  commit fork --- it calls
  \texttt{Fx.CreateOutput(reward, validationRewardsOwner)} and adds it as
  a fresh reward UTXO
  (\texttt{proposal\_tx\_executor.go:300--322},
  \texttt{AddUTXO}/\texttt{AddRewardUTXO}). \IMPL{}
\end{enumerate}

\noindent Under \texttt{XStakeCommitment} all three move: (1)~P's
\texttt{utxoDB} is dropped (P holds claims and eligibility data, not
spendable UTXOs); (2)~the principal lives in an \texttt{XStakeCommitment}
on X, not in \texttt{StakeOuts} on P; (3)~principal return becomes
\texttt{CompleteStake}/\texttt{AbortStake}/\texttt{FailReward} on X, and
reward issuance becomes the X-level reward-mint transition. This is the
work of Phase P4 (Section~\ref{sec:migration}); it is the single clearest
place current code contradicts the model, already surfaced for the
DEX-vs-stake symmetry in Section~\ref{sec:per-chain}.

\paragraph{The commit/abort fork is the structural obstacle.} Beyond
minting, P's reward path has a structural dependency the DEX path does
not: the reward UTXO is written to \texttt{onCommitState} \emph{and} the
refund UTXO is written to \emph{both} \texttt{onCommitState} and
\texttt{onAbortState} (\texttt{proposal\_tx\_executor.go:294--295},
\texttt{300--322}). P's settlement is thus entangled with the
\emph{proposal commit/abort fork} --- a P-Chain BFT-vote construct (C2,
Section~\ref{sec:rewardmodel}) with no analog on X.
\texttt{CompleteStake}/\texttt{FailReward} must reproduce the
commit/abort economics (reward only on commit, principal returned on
either) as X settlements \emph{conditioned on the branch P's BFT vote
finalized}, \emph{without} importing P's proposal-vote machinery onto X
and \emph{without} X ever trusting a P-supplied reward amount. Getting
that conditioning right --- a reward minted on X exactly when P committed
and never when P aborted, recomputed by X --- is a P4 design obligation
that the transient DEX \texttt{SettleTx} never faces; it is recorded as a
distinct item for Section~\ref{sec:open}.

\subsection{Rewards are new value: an X-level inflation primitive}
\label{sec:rewardinflation}

Every other XSS transition is net-zero on total supply: \texttt{LockTx},
\texttt{UnlockTx}, the transient \texttt{SettleTx}, and the no-reward
staking transitions (\texttt{CommitStake}, \texttt{AcceptStake},
\texttt{AbortStake}, \texttt{FailReward}, \texttt{RevokeOperator}) only
move value between the spendable, locked, and shielded pools
(Section~\ref{sec:safety}, S2). \textbf{Staking rewards are not.} A
reward is \emph{new} value of asset $a$, created at \texttt{CompleteStake},
and it therefore cannot be expressed as a net-zero \texttt{SettleTx}.

\CRYPTO{}\,/\,\DESIGN{} XSS requires a \emph{distinct X-level inflation
primitive} --- the reward-mint transition --- that is the \emph{only}
transition allowed to increase $\mathrm{total}(a)$, and only by the
amount X's schedule permits. Its obligations:

\begin{itemize}[leftmargin=2em,nosep]
\item It fires \emph{only} inside \texttt{CompleteStake}, gated by the
  branch-finality conditioning of C2 (X reads which block branch P
  finalized; reward on the commit branch, none on the abort branch). X
  recomputes the amount from \texttt{XRewardSchedule} --- P proposes the
  branch, X authorizes and sizes the inflation.
\item It is the \emph{sole} writer of new supply. With it, the global
  conservation invariant (Section~\ref{sec:safety}, S2) is restated:
  $\mathrm{total}(a)$ is constant \emph{except} at the reward-mint
  transition (which increases it by exactly the authorized reward) and at
  fee burn (which decreases it). \textbf{There is no slash-burn term} ---
  no staking transition decreases supply by confiscating principal. Every
  other transition is net-zero.
\item The amount must be bound by an X-enforced schedule, not taken on
  P's word --- the legacy bound already exists as the supply-cap
  invariant (\texttt{state\_changes.go:173--175}). The precise schedule,
  the branch-finality proof that ties a reward to a specific bond and
  uptime, and the monetary-policy bound are cryptographer- and
  economics-review items.
\end{itemize}

\noindent This connects to the open question already recorded in
Section~\ref{sec:open} (``P's mint authority \ldots\ reward minting
becomes an X-level inflation primitive gated by P's settlement, distinct
from the net-zero \texttt{SettleTx}''). \texttt{XStakeCommitment}
specifies \emph{where} that primitive fires (the \texttt{CompleteStake}
transition), \emph{what} conditions it (the C2 branch-finality proof), and
\emph{that supply changes only at reward-mint and fee-burn} --- never at a
slash; it does not assert the schedule or the proof are sound, which
remain gated on sign-off.

\subsection{The final layered model}
\label{sec:stake-layers}

The staking design is one instance of a clean layering, which we state
because the rest of the LP relies on it:

\begin{itemize}[leftmargin=2em,nosep]
\item \textbf{X} owns money: UTXOs, stake commitments, the reward-mint
  schedule, and the capability NFTs.
\item \textbf{P} owns the validator lifecycle: it proves completion,
  failure, and eligibility, and cannot mint, burn, or confiscate.
\item \textbf{D / B} own specialized execution: they produce receipts and
  proofs X settles.
\item \textbf{Z} proves shielded X-state transitions across the boundary.
\item \textbf{Misbehavior} costs an actor its reward, its capability, its
  route, or its operator status --- never another actor's or its own
  principal.
\end{itemize}
\noindent \textbf{The base protocol has no slashing.}

\subsection{Test plan for the X-staking build (design-first TDD)}
\label{sec:stake-tests}

These are the validation obligations for the staking-on-X migration; they
are \DESIGN{} targets for Phase P4 (Section~\ref{sec:migration}), written
now so the migration is test-driven, \emph{not} to be implemented before
P4. The slash-era cases are removed because the transition they tested no
longer exists.

\paragraph{Removed (the slash transition is gone).}
\begin{itemize}[leftmargin=2em,nosep]
\item \texttt{TestSlashStakeBurnsPrincipal}
\item \texttt{TestSlashRedirectsStake}
\item \texttt{TestSlashedStakeConservationWithBurn}
\end{itemize}

\paragraph{Added (the no-slash lattice).}
\begin{itemize}[leftmargin=2em,nosep]
\item \texttt{TestOfflineValidatorReceivesPrincipalNoReward} ---
  \texttt{FailReward} returns full principal, zero reward.
\item \texttt{TestInvalidLifecycleProofCannotConsumeStake} --- a malformed
  completion/failure proof cannot set \texttt{consumed}.
\item \texttt{TestPChainCannotMintRewardDirectly} --- no P transition
  creates a reward output; only X's reward-mint does.
\item \texttt{TestRewardMintBoundedByXSchedule} --- the minted reward
  equals \texttt{XRewardSchedule} and respects the supply cap.
\item \texttt{TestEquivocationEvidenceRevokesCapabilityNoBurn} ---
  recorded evidence revokes the NFT and burns nothing.
\item \texttt{TestOperatorNFTRevocationBlocksFuturePluginActivation} ---
  a revoked operator cannot re-activate the module.
\item \texttt{TestStakeCommitmentPrincipalAlwaysReturnedOnAbortOrFailReward}
  --- principal conservation on both no-reward terminals.
\end{itemize}

\noindent These map onto the validation checklist
(Section~\ref{sec:validation}): the principal-return cases extend P8
(commitment exclusivity / consume-once), the reward-bound case extends
the S2/P2 conservation accounting, and the NFT cases extend P5
(capability gate fails closed).
